\documentclass[11pt,letterpaper]{article}
% This source is intentionally restricted to pdfLaTeX.
\usepackage[T1]{fontenc}
\usepackage[utf8]{inputenc}
\usepackage{textcomp}
\usepackage{lmodern}
\usepackage[margin=0.82in,top=0.94in,bottom=0.91in,headheight=18pt]{geometry}
\usepackage{amsmath,amssymb}
\usepackage[table]{xcolor}
\definecolor{DeepNavy}{HTML}{14283F}
\definecolor{Teal}{HTML}{0A747C}
\definecolor{Gold}{HTML}{C78F36}
\definecolor{Slate}{HTML}{536477}
\definecolor{Pale}{HTML}{F4F7FA}
\definecolor{CodeEdge}{HTML}{D4DFE9}
\usepackage{microtype}
\usepackage{parskip}
\setlength{\parskip}{0.47em}
\usepackage{titlesec}
\titleformat{\section}{\normalfont\Large\bfseries\color{DeepNavy}}{\thesection}{0.62em}{}
\titlespacing*{\section}{0pt}{2.0ex plus .7ex minus .2ex}{0.75ex}
\titleformat{\subsection}{\normalfont\large\bfseries\color{Teal}}{\thesubsection}{0.6em}{}
\usepackage{longtable,booktabs,array,calc}
\renewcommand{\arraystretch}{1.22}
\AtBeginEnvironment{longtable}{\small}
\usepackage{etoolbox}
\makeatletter
\patchcmd\longtable{\par}{\if@noskipsec\mbox{}\fi\par}{}{}
\makeatother
\usepackage{listings}
\IfFileExists{upquote.sty}{\usepackage{upquote}}{}
\lstdefinestyle{suitecode}{
  basicstyle=\ttfamily\footnotesize,
  breaklines=true,breakatwhitespace=false,
  columns=fullflexible,keepspaces=true,showstringspaces=false,
  frame=single,rulecolor=\color{CodeEdge},
  backgroundcolor=\color{Pale},
  keywordstyle=\color{Teal}\bfseries,
  commentstyle=\color{Slate}\itshape,
  stringstyle=\color{DeepNavy},
  xleftmargin=5pt,xrightmargin=5pt,
  framexleftmargin=5pt,framexrightmargin=5pt,
  aboveskip=0.85em,belowskip=0.95em,
  tabsize=2
}
\lstset{style=suitecode}
\usepackage{enumitem}
\setlist[itemize]{leftmargin=1.6em,itemsep=2pt,topsep=3pt}
\setlist[enumerate]{leftmargin=1.8em,itemsep=2pt,topsep=3pt}
\providecommand{\tightlist}{\setlength{\itemsep}{0pt}\setlength{\parskip}{0pt}}
\newcommand{\passthrough}[1]{#1}
\usepackage{xurl}
\usepackage[colorlinks=true,linkcolor=Teal,urlcolor=Teal,citecolor=Teal,pdfborder={0 0 0}]{hyperref}
\urlstyle{same}
\usepackage{fancyhdr,lastpage}
\pagestyle{fancy}
\fancyhf{}
\fancyhead[L]{\footnotesize\sffamily\color{Slate} ECMA-262:2025 \textbullet{} Engineering Documentation}
\fancyhead[R]{\footnotesize\sffamily\color{Slate} 16th Edition}
\fancyfoot[L]{\footnotesize\sffamily\color{Slate} Independent technical companion}
\fancyfoot[R]{\footnotesize\sffamily\color{Slate}\thepage\ / \pageref{LastPage}}
\renewcommand{\headrulewidth}{0.35pt}
\renewcommand{\footrulewidth}{0pt}
\setlength{\emergencystretch}{3em}
\setcounter{tocdepth}{2}
\setcounter{secnumdepth}{2}
\begin{document}
\begin{titlepage}
\thispagestyle{empty}
\vspace*{2.8cm}
{\color{Gold}\rule{18mm}{3pt}}\par\vspace{18mm}
{\sffamily\Large\bfseries\color{Teal} SOFTWARE ENGINEERING DOCUMENTATION SUITE}\par
\vspace{8mm}
{\sffamily\huge\bfseries\color{DeepNavy} ECMAScript 2025 }\par
\vspace{4mm}
{\sffamily\fontsize{24}{29}\selectfont\bfseries\color{DeepNavy} Specification Study Guide }\par
\vspace{10mm}
{\large\color{Slate} Clause-based interpretation and assessments }\par
\vspace{16mm}
{\color{Teal}\rule{\textwidth}{0.8pt}}\par
\vfill
{\sffamily\color{Slate}\large ECMA-262:2025 \textbullet{} 16th Edition \textbullet{} June 2025}\par
\vspace{3mm}
{\sffamily\color{Slate} October 2026 \textbullet{} Document 2 of 8}\par
\vspace{3mm}
{\small\color{Slate} Original engineering companion. The uploaded ECMA-262:2025 specification remains the authoritative source.}\par
\end{titlepage}
\pagenumbering{arabic}
\tableofcontents
\clearpage
\section{Purpose and study discipline}\label{purpose-and-study-discipline}

This clause-oriented study guide teaches how to \emph{interpret} ECMA-262:2025 rather than memorize isolated API names. A learner should be able to resolve a surprising JavaScript expression by following grammar, static semantics, abstract operations, and execution algorithms. All clause numbers follow the uploaded June 2025 edition; an online living edition may be renumbered.

\textbf{Evidence legend:} \textbf{{[}SPEC{]}} = summarized standard semantics; \textbf{{[}HOST{]}} = integration behavior not fully determined by ECMA-262; \textbf{{[}GUIDE{]}} = study technique or recommended application behavior. Clause 2's legacy and normative-optional designations are important: do not reinterpret all informative annex material as mandatory.

\section{Learning sequence and source map}\label{learning-sequence-and-source-map}

\begin{longtable}[]{@{}
  >{\raggedright\arraybackslash}p{(\columnwidth - 4\tabcolsep) * \real{0.3333}}
  >{\raggedright\arraybackslash}p{(\columnwidth - 4\tabcolsep) * \real{0.3333}}
  >{\raggedright\arraybackslash}p{(\columnwidth - 4\tabcolsep) * \real{0.3333}}@{}}
\toprule\noalign{}
\begin{minipage}[b]{\linewidth}\raggedright
Study block
\end{minipage} & \begin{minipage}[b]{\linewidth}\raggedright
Read
\end{minipage} & \begin{minipage}[b]{\linewidth}\raggedright
Evidence of mastery
\end{minipage} \\
\midrule\noalign{}
\endhead
\bottomrule\noalign{}
\endlastfoot
1. Specification literacy & 1--5, 17, Annexes & explain normative/informative boundaries \\
2. Values and coercion & 6--7 & trace \passthrough{\lstinline!ToPrimitive!}, equality, Numeric types \\
3. Parsing and binding & 8, 11--14 & identify early errors, TDZ, ASI \\
4. Execution & 9--10 & describe realm, context, records, proxy traps \\
5. Functions and modules & 15--16 & trace closure, async, module evaluation \\
6. Standard library & 18--28 & navigate intrinsic methods and invariants \\
7. Shared memory & 25, 29 & distinguish atomic and non-atomic access \\
8. Edition deltas & 16, 21--25, 27 & demonstrate 2025 additions \\
\end{longtable}

\section{Block 1: Specification literacy (1--5, 17)}\label{block-1-specification-literacy-15-17}

\textbf{{[}SPEC{]}} Scope defines the language; conformance constrains implementations. Grammar distinguishes lexical and syntactic productions. Early errors are not ordinary runtime exceptions that can be caught after code starts running. \passthrough{\lstinline!?!} propagates an abrupt completion; \passthrough{\lstinline"!"} indicates that the specification expects an operation's completion to be normal in that context. A specification \emph{Reference Record} is an internal abstraction for resolving variable/property references, not an application-level object.

\textbf{Worked trace.} \passthrough{\lstinline!obj.x!} can create a reference, resolve a base, and invoke \passthrough{\lstinline![[Get]]!}; a getter or proxy trap can execute arbitrary user code. The string \passthrough{\lstinline!obj.x!} alone does not reveal whether accessing the property is pure.

\textbf{Checkpoint:} Why would conformance permits host extensions but forbid extensions listed in 17.1? How can a normative-optional clause differ from informative text? When is browser Annex B required?

\section{Block 2: Values, conversions, and comparisons (6--7)}\label{block-2-values-conversions-and-comparisons-67}

\textbf{{[}SPEC{]}} BigInt and Number are separate numeric types; \passthrough{\lstinline!ToNumeric!} does not authorize combining them with every arithmetic operator. SameValue, SameValueZero, and strict equality differ, particularly around \passthrough{\lstinline!NaN!} and signed zero. String comparison uses UTF-16 code-unit behavior in language-defined contexts; do not conflate with natural-language collation (ECMA-402 is separate).

\begin{lstlisting}
console.assert(0n == 0);         // true with abstract equality
console.assert(0n !== 0);        // distinct language types
console.assert(Object.is(-0, 0) === false);
console.assert(new Map([[NaN, 'x']]).get(NaN) === 'x');
\end{lstlisting}

\textbf{Checkpoint:} Given \passthrough{\lstinline![] == false!}, document the abstract conversion steps and contrast with \passthrough{\lstinline![] === false!}. For \passthrough{\lstinline!1n + 1!}, identify whether it produces a number or throws and why.

\section{Block 3: Grammar, expressions, and declarations (8, 11--14)}\label{block-3-grammar-expressions-and-declarations-8-1114}

\textbf{{[}SPEC{]}} Lexical grammar is context-sensitive in specific places; there are multiple lexical goal symbols. Statements and declarations have static and runtime semantics. \passthrough{\lstinline!let!}/\passthrough{\lstinline!const!} bindings are created but not initialized immediately; accesses in the temporal dead zone throw. Modules are always strict-mode code. ASI is specified at 12.10 and can alter observable \passthrough{\lstinline!return!} behavior.

\begin{lstlisting}
// The function body is not evaluated until called.
function readBeforeInit() {
  return item;
  let item = 3;
}
try { readBeforeInit(); } catch (e) {
  console.assert(e instanceof ReferenceError);
}
\end{lstlisting}

\textbf{Checkpoint:} Describe why \passthrough{\lstinline!return!} followed by a newline and object literal returns \passthrough{\lstinline!undefined!}. Explain difference between a syntax error and a runtime \passthrough{\lstinline!ReferenceError!}.

\section{Block 4: Realms, contexts, and object semantics (9--10)}\label{block-4-realms-contexts-and-object-semantics-910}

\textbf{{[}SPEC{]}} Each realm has its own intrinsics. An execution context records a running evaluation state and associated environment structures. Ordinary objects implement essential internal methods, while exotic objects may require special logic. A proxy can trap operations subject to invariants, such as those around nonconfigurable properties.

\begin{lstlisting}
const target = Object.defineProperty({}, 'id', {
  value: 7, writable: false, configurable: false
});
const invalid = new Proxy(target, { get() { return 8; } });
try { void invalid.id; } catch (e) {
  console.assert(e instanceof TypeError);
}
\end{lstlisting}

\textbf{Checkpoint:} Explain why two arrays from different realms may fail an \passthrough{\lstinline!instanceof!} check, yet pass \passthrough{\lstinline!Array.isArray!}. Distinguish a proxy invariant from application authorization policy.

\section{Block 5: Functions, classes, and asynchronous modules (15--16)}\label{block-5-functions-classes-and-asynchronous-modules-1516}

\textbf{{[}SPEC{]}} Arrow functions capture \passthrough{\lstinline!this!} lexically; class bodies and modules obey strict-mode semantics; generators suspend and resume; asynchronous functions return promises. Imports are live bindings, not copied snapshots. Module loading and resolution are host-defined integration points, including attributes accepted by a host.

\begin{lstlisting}
function counter() {
  let n = 0;
  return () => ++n;
}
const next = counter();
console.assert(next() === 1 && next() === 2);
\end{lstlisting}

\textbf{Checkpoint:} Differentiate lexical environments, private environments, and module environment records. Where is the dependency graph linked versus evaluated? Which part is host policy?

\section{Block 6: Intrinsics, collections, iteration, and promises (18--28)}\label{block-6-intrinsics-collections-iteration-and-promises-1828}

\textbf{{[}SPEC{]}} Built-in objects have specified attributes, receiver validations, iteration contracts, and error behavior. \passthrough{\lstinline!Array.prototype.map!} may encounter holes; \passthrough{\lstinline!Map!} maintains key insertion semantics; typed-array methods operate on numeric storage; weak references deliberately do not expose predictable garbage-collection timing. Async promise reactions are jobs, not immediate invocations inline with \passthrough{\lstinline!.then!}.

\begin{lstlisting}
const log = [];
Promise.resolve().then(() => log.push('reaction'));
log.push('current job');
// Check ordering at a later host-controlled job checkpoint.
\end{lstlisting}

\textbf{Checkpoint:} Why is \passthrough{\lstinline!WeakRef!} inappropriate for application-critical bookkeeping? What happens if iterator \passthrough{\lstinline!.next()!} throws during consumption? How do array holes differ from present \passthrough{\lstinline!undefined!} values?

\section{Block 7: Shared memory and observable concurrency (25, 29)}\label{block-7-shared-memory-and-observable-concurrency-25-29}

\textbf{{[}SPEC{]}} \passthrough{\lstinline!SharedArrayBuffer!}, \passthrough{\lstinline!Atomics!}, agents, and the memory model are separate from promise jobs. Shared reads/writes, ordering relations, and data-race concepts matter when multiple agents execute. \passthrough{\lstinline!Atomics!} provides memory-model semantics for supported typed arrays, but cannot turn arbitrary application state into a synchronized transaction.

\begin{lstlisting}
// Shared buffer can be shared between agents by host mechanisms.
const buffer = new SharedArrayBuffer(4);
const words = new Int32Array(buffer);
Atomics.store(words, 0, 5);
console.assert(Atomics.load(words, 0) === 5);
\end{lstlisting}

\textbf{Checkpoint:} Explain why demonstration within one agent does not validate cross-agent ordering. Identify the host mechanism needed to start workers.

\section{Block 8: 2025 additions and comparison with 2024}\label{block-8-2025-additions-and-comparison-with-2024}

\textbf{{[}SPEC{]}} The edition preface lists Iterator helpers, Set methods, JSON modules with import attributes, \passthrough{\lstinline!RegExp.escape!}, inline regexp modifiers, \passthrough{\lstinline!Promise.try!}, and \passthrough{\lstinline!Float16Array!} with associated APIs. Do not falsely classify 2024 features (\passthrough{\lstinline!Promise.withResolvers!}, resizable buffers, \passthrough{\lstinline!Atomics.waitAsync!}) as introduced in 2025.

\begin{longtable}[]{@{}
  >{\raggedright\arraybackslash}p{(\columnwidth - 4\tabcolsep) * \real{0.3333}}
  >{\raggedright\arraybackslash}p{(\columnwidth - 4\tabcolsep) * \real{0.3333}}
  >{\raggedright\arraybackslash}p{(\columnwidth - 4\tabcolsep) * \real{0.3333}}@{}}
\toprule\noalign{}
\begin{minipage}[b]{\linewidth}\raggedright
2025 change
\end{minipage} & \begin{minipage}[b]{\linewidth}\raggedright
Investigation task
\end{minipage} & \begin{minipage}[b]{\linewidth}\raggedright
Target clause
\end{minipage} \\
\midrule\noalign{}
\endhead
\bottomrule\noalign{}
\endlastfoot
Iterator helpers & identify laziness and terminal consumption & 27.1 \\
Set algebra & check comparator and result insertion behavior & 24.2 \\
JSON modules & separate grammar from host loading & 16.2 \\
RegExp escaping & explain leading ASCII alphanumeric escaping & 22.2.5.1 \\
Inline modifiers & construct scoped flag-sensitive pattern & 22.2 \\
Promise.try & compare thrown exception vs returned thenable & 27.2.4.8 \\
Float16 storage & reason about rounding/representability & 21, 23, 25 \\
\end{longtable}

\section{Guided specification reading: trace a setter}\label{guided-specification-reading-trace-a-setter}

\begin{enumerate}
\def\labelenumi{\arabic{enumi}.}
\tightlist
\item
  Start with an assignment expression in clause 13.
\item
  Determine the reference being assigned to and any left-side evaluation effects.
\item
  Follow \passthrough{\lstinline!PutValue!} / \passthrough{\lstinline!Set!} to object internal method dispatch.
\item
  Determine whether the own property is data or accessor, whether a prototype setter applies, and whether a proxy trap intervenes.
\item
  Record completion behavior (normal or throw), receiver identity, and strict-mode differences.
\item
  Write tests for accessor success, receiver mismatch, non-writable property, and proxy invariant failure.
\end{enumerate}

\textbf{{[}GUIDE{]}} A worked trace is complete only when each transition is attributable to a specific algorithm or host hook rather than an intuition about an engine's implementation.

\section{Assessment bank}\label{assessment-bank}

\begin{enumerate}
\def\labelenumi{\arabic{enumi}.}
\tightlist
\item
  Compare the eight language types with specification-only types and give an example of each category.
\item
  Contrast SameValue, SameValueZero, strict equality, and abstract equality with counterexamples.
\item
  Define early errors, static semantics, runtime semantics, and abrupt completion.
\item
  Explain realms, global environments, and per-realm intrinsics.
\item
  Trace \passthrough{\lstinline!this!} in strict functions, methods, arrows, and class initializers.
\item
  Discuss proxy invariants and how they constrain observable traps.
\item
  Explain dynamic import vs static import dependency and the role of the host.
\item
  Explain why Promise reactions run asynchronously.
\item
  Identify two cases in which iteration invokes user-defined code.
\item
  Design a test matrix for regexp escaping, half-precision rounding, or set intersection.
\end{enumerate}

\textbf{Mastery rubric:} 0 = unsupported claim; 1 = recognizes the right subject; 2 = correct observable answer; 3 = traces the applicable clause and edge cases; 4 = can write a discriminating test and document host assumptions. Recommended pass threshold: at least level 3 on every block.

\section{Primary reference}\label{primary-reference}

Ecma International, \emph{ECMA-262, ECMAScript 2025 Language Specification}, 16th edition, June 2025 (user-supplied PDF). Secondary current living navigation: https://tc39.es/ecma262/ ; do not silently substitute live-edition clause numbers.

\section{Appendix: top-level clause-by-clause examination index}\label{appendix-top-level-clause-by-clause-examination-index}

This catalog supplies a prompt for \textbf{each of the 29 numbered top-level clauses} in the source edition; each row is a learning aid rather than a substitute for individual subclause algorithms.

\begin{longtable}[]{@{}
  >{\raggedright\arraybackslash}p{(\columnwidth - 4\tabcolsep) * \real{0.3333}}
  >{\raggedright\arraybackslash}p{(\columnwidth - 4\tabcolsep) * \real{0.3333}}
  >{\raggedright\arraybackslash}p{(\columnwidth - 4\tabcolsep) * \real{0.3333}}@{}}
\toprule\noalign{}
\begin{minipage}[b]{\linewidth}\raggedright
Clause
\end{minipage} & \begin{minipage}[b]{\linewidth}\raggedright
Standard topic
\end{minipage} & \begin{minipage}[b]{\linewidth}\raggedright
Spec-reading challenge
\end{minipage} \\
\midrule\noalign{}
\endhead
\bottomrule\noalign{}
\endlastfoot
1 & Scope & State exactly which language the edition defines. \\
2 & Conformance & Classify a legacy versus normative-optional clause. \\
3 & Normative References & Identify which external standards influence text and numbers. \\
4 & Overview & Explain the distinction between ECMAScript and a host. \\
5 & Notational Conventions & Explain grammar parameters and completion shorthand. \\
6 & Data Types and Values & Separate language values and specification Records. \\
7 & Abstract Operations & Trace \passthrough{\lstinline!ToPrimitive!} through an overloaded operator. \\
8 & Syntax-Directed Operations & Locate scope and bound-name analysis. \\
9 & Executable Code and Execution Contexts & Draw contexts, realms, environments, jobs, agents. \\
10 & Ordinary and Exotic Objects Behaviours & Trace \passthrough{\lstinline![[Get]]!} through a proxy invariant. \\
11 & Source Text & Identify source-code categories and text constraints. \\
12 & Lexical Grammar & Determine ASI result and regular-expression lexical goal. \\
13 & Expressions & Evaluate reference and operator order precisely. \\
14 & Statements and Declarations & Distinguish early error and runtime control completion. \\
15 & Functions and Classes & Compare lexical \passthrough{\lstinline!this!}, generators, private names. \\
16 & Scripts and Modules & Distinguish parse, link, host load, and evaluate. \\
17 & Error Handling and Language Extensions & Explain forbidden extensions. \\
18 & Standard Built-in Objects & Interpret common built-in conventions and attributes. \\
19 & The Global Object & Identify property categories and global constructors. \\
20 & Fundamental Objects & Verify descriptor and function-object behavior. \\
21 & Numbers and Dates & Compare Number, BigInt, Float16 rounding, Date. \\
22 & Text Processing & Trace String/RegExp algorithms and escaping. \\
23 & Indexed Collections & Distinguish array holes, TypedArray buffer semantics. \\
24 & Keyed Collections & Trace Map/Set SameValueZero and Set algebra. \\
25 & Structured Data & Analyze ArrayBuffer, DataView, Atomics, JSON. \\
26 & Managing Memory & Explain weak reachability and cleanup nondeterminism. \\
27 & Control Abstraction Objects & Trace iterator, promise, generator, async semantics. \\
28 & Reflection & Identify Reflect/Proxy method invariants. \\
29 & Memory Model & Describe agent event relations and data-race freedom. \\
\end{longtable}

\textbf{Annex reading order:} A for informative grammar consolidation; B for web-browser normative compatibility; C for informative strict-mode synthesis; D for host layering points; E--F for informative historical corrections and incompatible changes. Cite an annex's status whenever presenting it as evidence for a required implementation behavior.

\section{Appendix: worked investigation protocols}\label{appendix-worked-investigation-protocols}

\textbf{Algorithm comparison protocol:} Choose two similar operations (for example, SameValue and SameValueZero). Record parameters, return type, branch points, exceptions, users of each operation, and at least two inputs on which results differ. State when the operation is observable directly versus only through a built-in that delegates to it.

\textbf{Edition-change protocol:} Start with the 2025 foreword, locate the newly standardized method's normative subclause, inspect receiver validation and completion behavior, and write one success, one error, and one side-effect test. Then determine whether the target runtime supports the feature; avoid treating 2025 publication as proof of immediate universal deployment.

\textbf{Host-boundary protocol:} For a module import or promise rejection observation, identify which events are language-defined, which host hooks are invoked, and what the target host documentation must specify. Do not claim that host integrations are entirely unconstrained; rather, distinguish the obligations ECMA-262 places on hooks from the policy those hooks are permitted to choose.

\end{document}
